% Pollis drill-down: Plotting, charts and maps drawn by Pollis: why they are drawn in JavaScript, the libraries for
% charts, maps, street-level imagery and 3D, the Julia packages that describe them, where each appears, and what is
% still to come (21 November 2026). For users who want the detail, and for whoever maintains Pollis.
% Compile with: latexmk -pdf workflow_plotting.tex
% Shared look for the one-page Pollis workflow handouts (same fonts and colours as the brochure).
\documentclass[10pt,a4paper]{article}

\usepackage[a4paper,left=17mm,right=17mm,top=14mm,bottom=16mm,footskip=8mm]{geometry}
\usepackage[T1]{fontenc}
\usepackage[default,semibold]{sourcesanspro}
\usepackage{microtype}
\usepackage{graphicx}
\usepackage[table]{xcolor}
\usepackage{tikz}
\usepackage[most]{tcolorbox}
\usepackage{fontawesome5}
\usepackage{enumitem}
\usepackage{fancyhdr}
\usepackage[hidelinks]{hyperref}

\graphicspath{{images/}}

% ---------------------------------------------------------------- colours (as in the brochure)
\definecolor{navy}{HTML}{1C3A5E}
\definecolor{pblue}{HTML}{3B78B5}
\definecolor{amber}{HTML}{F2B33D}
\definecolor{deepamber}{HTML}{D9902A}
\definecolor{pale}{HTML}{EEF3F9}
\definecolor{ink}{HTML}{2A2F36}
\definecolor{logonavy}{HTML}{173D63}
\definecolor{logodark}{HTML}{C98A1E}
\definecolor{logoamber}{HTML}{F2B446}
\definecolor{logocream}{HTML}{FEEFD2}
\definecolor{soft}{HTML}{5E6773}
\color{ink}

\setlength{\parindent}{0pt}
\setlength{\parskip}{3pt}
\setlist{nosep,leftmargin=*,itemsep=2pt}

\pagestyle{fancy}
\fancyhf{}
\renewcommand{\headrulewidth}{0pt}
\fancyfoot[L]{\footnotesize\color{soft}\textbf{Pollis}\enspace Bright Insights through Exploration, Modelling, Simulation and Optimisation}
\fancyfoot[R]{\footnotesize\color{soft}github.com/saragga/pollis/releases/latest}

% The Pollis logo, redrawn so it prints sharply at any size
\newcommand{\pollislogo}[1]{%
  \begin{tikzpicture}[x=#1/512,y=-#1/512,baseline=(current bounding box.center)]
    \fill[logonavy] (0,0) rectangle (512,512);
    \fill[logodark] (36,422) -- (104,256) -- (172,422) -- cycle;
    \fill[logoamber] (104,422) -- (203,180) -- (302,422) -- cycle;
    \fill[logocream] (203,422) -- (339,89) -- (476,422) -- cycle;
  \end{tikzpicture}}

% Top of the page: logo, name, kind of document (Workflow unless given, e.g. [Drill-Down]) and its subject;
% then the question and a one-line promise
\newcommand{\workflowhead}[4][Workflow]{%
  \thispagestyle{fancy}%
  \begin{minipage}[c]{13mm}\pollislogo{11mm}\end{minipage}%
  \begin{minipage}[c]{\dimexpr\linewidth-13mm\relax}
    {\color{navy}\Large\bfseries Pollis}\enspace{\color{soft}#1}\hfill
    {\color{deepamber}\small\bfseries\MakeUppercase{#2}}
  \end{minipage}\par\vspace{7mm}
  {\color{navy}\fontsize{22}{26}\selectfont\bfseries #3}\par\vspace{4pt}
  {\color{soft}\large #4}\par\vspace{5mm}}

% A menu path: \menu{Explore}{Access Data Libraries}... is written as \path{Explore \gt Access Data Libraries}
\newcommand{\gt}{\,{\color{deepamber}\bfseries ›}\,}
\newcommand{\menupath}[1]{{\footnotesize\bfseries\color{navy}#1}}

% A screenshot if images/<file> exists, otherwise an orange placeholder that says what to capture
\newlength{\shotheight}\setlength{\shotheight}{43mm}
\newcommand{\stepshot}[2]{%
  \IfFileExists{images/#1}{%
    \tcbox[enhanced, colback=black, colframe=black!70, boxrule=0.4pt, arc=2pt,
      left=0pt,right=0pt,top=0pt,bottom=0pt, boxsep=0pt, halign=center]{%
      \includegraphics[width=\dimexpr\linewidth-1pt\relax,height=\shotheight,keepaspectratio]{#1}}}{%
    \begin{tcolorbox}[enhanced, colback=white, colframe=orange!80, boxrule=0.6pt, arc=2pt,
      height=\shotheight, valign=center, halign=center, left=6pt, right=6pt]
      {\color{orange}\Large\faCamera}\par\smallskip
      {\color{orange}\footnotesize\bfseries Screenshot: #2}\par\smallskip
      {\color{orange!80!black}\scriptsize images/#1}
    \end{tcolorbox}}}

% One step: number, title, menu path, image file, what to capture, what it means
\newcommand{\step}[6]{%
  \begin{tcolorbox}
    {\color{deepamber}\Large\bfseries #1}\enspace{\color{navy}\large\bfseries #2}\par\vspace{1pt}
    \menupath{#3}\par\vspace{4pt}
    \stepshot{#4}{#5}\par\vspace{3pt}
    {\small #6}
  \end{tcolorbox}}

\newenvironment{steps}{%
  \begin{tcbraster}[raster columns=2, raster equal height=rows, raster column skip=4mm,
    raster row skip=4mm, enhanced, colback=pale, colframe=pale, boxrule=0pt, sharp corners,
    left=6pt,right=6pt,top=6pt,bottom=6pt, fontupper=\raggedright]}{\end{tcbraster}}

% Boxes for the bottom of the page
\newtcolorbox{pbox}[2][]{enhanced, colback=pale, colframe=amber, boxrule=0pt,
  leftrule=3pt, sharp corners, left=7pt, right=7pt, top=6pt, bottom=6pt,
  fontupper=\small\raggedright, before upper={{\color{navy}\normalsize\bfseries #2}\par\smallskip}, #1}
\newtcolorbox{nbox}[1][]{enhanced, colback=navy, colframe=navy, boxrule=0pt,
  sharp corners, left=8pt, right=8pt, top=6pt, bottom=6pt, coltext=white,
  fontupper=\small\raggedright, #1}

% A Next Steps item inside the bottom boxes
\newcommand{\nsitem}[2]{\item[\color{deepamber}\faChevronRight] \textbf{#1}: #2}

% Citations: \dcite{key}{Breiman (2001)} links to its entry \dref{key}{...} in the References, keyed as in pollis.bib
\newcommand{\dcite}[2]{\hyperlink{ref:#1}{\color{pblue}#2}}
\newcommand{\dref}[2]{\raisebox{1.1em}[0pt][0pt]{\hypertarget{ref:#1}{}}#2\par}
\newcommand{\ddoi}[1]{\href{https://doi.org/#1}{\color{pblue}doi:#1}}
\newcommand{\references}{\section*{References}\addcontentsline{toc}{section}{References}\footnotesize\raggedright\setlength{\parskip}{3pt}}

% The page of each heading, for Ask Pollis
% Pollis: the page of each heading, for Ask Pollis. A document that inputs this file, at the end of its preamble, writes
% <jobname>.labels beside its PDF, a line page|level|heading for each heading, the page counted from the first, the
% cover, as 1. compile_documents.jl copies it to ~/.pollis/documents, and ask_index.jl cuts the document into passages
% under its headings with it. \asklabel{1}{Title} marks a topic or a page, \asklabel{2}{Title} a heading inside it; the
% sections and subsections of a document that has them are marked as they go into its contents. A package vignette
% inputs it before its \documentclass (compile_vignette.jl), so the contents are taken over when the document begins,
% after the packages have taken them over. Read a second time, as by a copy of a document for Edit the LaTeX Source,
% which inputs it as resources/ask_labels when Pollis has put it in already (document_source.jl), it does nothing.
\ifdefined\asklabel\endinput\fi
\makeatletter
\newwrite\pollis@labels
\immediate\openout\pollis@labels=\jobname.labels
\newcommand{\asklabel}[2]{\write\pollis@labels{\the\ReadonlyShipoutCounter|#1|\unexpanded{#2}}}
\def\pollis@toc{toc}
\newcommand{\pollis@contentsline}[3]{%
  \def\pollis@file{#1}%
  \ifx\pollis@file\pollis@toc\asklabel{#2}{#3}\fi
  \pollis@addcontentsline{#1}{#2}{#3}}
% The headings of the New Methods style of the vignettes, \nmsection and \nmsub, which go into no contents
\def\pollis@asksection#1{\pollis@nmsection{#1}\asklabel{1}{#1}}
\def\pollis@asksub#1{\pollis@nmsub{#1}\asklabel{2}{#1}}
\AddToHook{begindocument/end}{%
  \let\pollis@addcontentsline\addcontentsline\let\addcontentsline\pollis@contentsline
  \ifdefined\nmsection\let\pollis@nmsection\nmsection\let\nmsection\pollis@asksection\fi
  \ifdefined\nmsub\let\pollis@nmsub\nmsub\let\nmsub\pollis@asksub\fi}
\makeatother



\usepackage{titlesec}
\usepackage{tabularx}
\usepackage{booktabs}
\usepackage{listings}

\fancyhf{}
\fancyfoot[L]{\footnotesize\color{soft}\textbf{Pollis}\enspace Drill-down: Plotting}
\fancyfoot[R]{\footnotesize\color{soft}\thepage}

\titleformat{\section}{\color{navy}\large\bfseries}{\color{deepamber}\thesection}{0.6em}{}
\titleformat{\subsection}{\color{navy}\normalsize\bfseries}{\color{deepamber}\thesubsection}{0.5em}{}
\titlespacing*{\section}{0pt}{12pt}{4pt}
\titlespacing*{\subsection}{0pt}{8pt}{2pt}
\setlength{\parskip}{4pt}

\lstset{basicstyle=\ttfamily\footnotesize, columns=fullflexible, keepspaces=true, breaklines=true,
  backgroundcolor=\color{pale}, frame=leftline, rulecolor=\color{amber}, framerule=2pt, xleftmargin=6pt,
  framexleftmargin=6pt, aboveskip=4pt, belowskip=4pt, upquote=true, tabsize=2}

\usepackage{xurl}
\newcommand{\file}[1]{{\small\nolinkurl{#1}}}
\newcolumntype{L}{>{\raggedright\arraybackslash}X}
\newcommand{\ui}[1]{\textbf{#1}}
% Code with spaces in it, which \file would run together
\newcommand{\code}[1]{\texttt{\small #1}}

\begin{document}

\workflowhead[Drill-Down]{Plotting}
  {Charts and maps drawn by Pollis}
  {Julia computes and describes; a JavaScript library in Pollis draws. Which library draws what, which Julia package
   describes it, where each chart and map appears, and what is still to come.}

\begin{nbox}
{\color{amber}\bfseries Drill-down.}\enspace
The charts are drawn by JavaScript in Pollis's own tab, so they appear without waiting for a Julia plotting package to
compile. ECharts covers the everyday charts, MapLibre covers the maps, and specialised JavaScript apps cover 3D and
city models.

This drill-down is for anyone who wants to know how a chart made in Julia reaches the screen, and for whoever
maintains Pollis. It describes the state of 21 November 2026, early in the work: the direction is set, and
\S\ref{sec:next} lists what is still to come.
\end{nbox}

\tableofcontents

\section{Why Pollis draws charts in JavaScript}\label{sec:why}
A Julia plotting package, such as Plots.jl, \dcite{CSRBB2023a}{Christ, Schwabeneder, Rackauckas, Borregaard and Breloff
(2023)}, or Makie.jl, \dcite{DK2021a}{Danisch and Krumbiegel (2021)}, draws the chart in Julia, and has to be compiled
before its first chart: on a computer where it has not been used, the first chart can take minutes. Pollis splits the
work in two.
\begin{itemize}
  \item \textbf{Julia computes and describes.} The code of a panel computes the numbers and writes what to draw: the
    kind of chart, its data, axes, colours and labels, as JSON, a short text.
  \item \textbf{JavaScript draws.} A library that Pollis ships, already loaded, reads the description and draws the
    chart in the tab, the notebook cell or the page.
\end{itemize}
The charts drawn this way are also interactive: a value shows when the pointer is over it, a series hides when its
name in the legend is clicked, and a chart that asks for them has zoom and a button that saves it as an image. Each kind of picture has the library
made for it (\S\ref{sec:libraries}).

\section{Three kinds of picture, their libraries}\label{sec:libraries}
\begin{center}\small
\renewcommand{\arraystretch}{1.25}
\begin{tabularx}{\linewidth}{@{}p{0.2\linewidth}p{0.3\linewidth}L@{}}
\toprule
\textbf{Picture} & \textbf{Library} & \textbf{Described in Julia by}\\
\midrule
Charts & Apache ECharts 6.0.0, \dcite{LMSSZ...2018a}{Li and 7 coauthors (2018)} & ECharts.jl for now; PollisECharts.jl
  (\S\ref{sec:charts})\\
Maps & MapLibre GL JS 5.24.0, \dcite{MLC2026a}{MapLibre Contributors (2026)}, on the vector tiles of OpenFreeMap,
  \dcite{Ero2026a}{Ero (2026)} & \file{Pollis.webmap} and \file{Pollis.layer}, of Pollis.jl (\S\ref{sec:maps})\\
Street-level imagery & MapillaryJS 4.1.2, \dcite{Map2023a}{Mapillary (2023)} & \file{Pollis.streetview}
  (\S\ref{sec:beyond})\\
Photorealistic 3D cities & The Maps JavaScript API, \dcite{Goo2026a}{Google (2026a)}, with its 3D maps,
  \dcite{Goo2026b}{Google (2026b)} & The Google Maps panel of Smart Regions (\S\ref{sec:beyond})\\
\bottomrule
\end{tabularx}
\end{center}
ECharts, MapLibre GL JS and MapillaryJS are bundled with Pollis, in the folder \file{libs} of the Julia extension, so
they draw without the internet; the base maps of MapLibre and the Google maps need it.

\section{Charts: ECharts}\label{sec:charts}
ECharts describes a chart as one JSON object, its \emph{option}: the series to draw and how, the axes, the legend,
the tooltip. It draws line, bar and scatter charts, box plots, heat maps, candlesticks, Sankey diagrams, trees,
graphs, calendars, radar and parallel coordinates, among others, on a canvas or as SVG.

\subsection{PollisECharts.jl}
The package that Pollis's charts are to be written with is \textbf{PollisECharts.jl}, a light package of Pollis's
own.
\begin{itemize}
  \item It depends only on JSON and Tables.jl: a table goes in, the option of the chart comes out.
  \item The chart is shown as its option alone, under a type of its own, as the maps of \file{Pollis.webmap} are
    (\S\ref{sec:maps}). Pollis draws it with the one copy of ECharts it ships, in the \ui{Plots} tab, in notebook cells
    and on pages, so a chart carries no copy of the library.
  \item It installs and loads in seconds, with no web server and no Node.js.
\end{itemize}
PollisECharts.jl is not written yet.

\subsection{ECharts.jl, for now}
Until then, the panels draw their charts with ECharts.jl 0.9, \dcite{Zwi2026a}{Zwitch (2026)}, a Julia package for
ECharts 6 with a function for each kind of chart (\file{line}, \file{bar}, \file{scatter}, \file{box},
\file{heatmap}, \file{sankey}, \file{candlestick} and many more). It serves the charts of Pollis well, but it is
heavier than Pollis needs, which is why it is a short-term choice:
\begin{itemize}
  \item it brings in HTTP.jl, for live charts that a small web server updates; NodeJS.jl, a copy of Node.js, to write
    a chart as an SVG file; and KernelDensity.jl, with FFTW, for density charts;
  \item the HTML of each chart carries its own copy of ECharts, about 1.1~MB. In the \ui{Plots} tab, Pollis takes the
    copy out and loads its own, of the same version, once; in notebook and Pluto cells, each chart keeps its copy.
\end{itemize}
On 21 November 2026, 27 panels of Pollis and its toolboxes draw with ECharts.jl.

\section{Maps: MapLibre}\label{sec:maps}
A map is written with functions of Pollis.jl, the module of Pollis's own Julia session.
\begin{lstlisting}
m = Pollis.webmap(Pollis.layer("districts.geojson"; name = "Districts"); style = "positron", title = "Districts")
\end{lstlisting}
\begin{itemize}
  \item \code{Pollis.webmap(layers...; style, center, zoom, title)}: a map with layers on an OpenStreetMap base map,
    drawn in the order given.
  \item \code{Pollis.layer(geojson; name, color, column, opacity)}: GeoJSON features, from a file or a text, in one
    colour or coloured by a column;
    \file{Pollis.tiles}: a layer of raster tiles, such as an image or a model's output.
  \item \file{Pollis.streetview}: a street-level viewer of Mapillary, alone or beside a map.
\end{itemize}
The map is shown as JSON of the type \file{application/vnd.pollis.map+json}, which Pollis draws with MapLibre GL JS.
On 21 November 2026, 31 panels draw maps this way: 21 in Smart Regions and 10 in Spatial Econometrics.

\section{Where a chart or a map appears}\label{sec:where}
\begin{center}\small
\renewcommand{\arraystretch}{1.25}
\begin{tabularx}{\linewidth}{@{}p{0.2\linewidth}LL@{}}
\toprule
\textbf{Code run in} & \textbf{A chart} & \textbf{A map}\\
\midrule
The Julia REPL, as with \ui{Julia REPL} below a panel's code & The \ui{Plots} tab, which keeps every chart: \ui{Plots: Show
  Previous Plot}, \ui{Plots: Save Plot} and the other commands of the tab, and the \ui{Plot Navigator} & A page in
  Pollis's browser, named after the map's title: a map shown again with the same title replaces it, others open
  beside it\\
A notebook, as with \ui{Notebook} & The cell's output & The cell's output, drawn by the renderer \ui{Pollis Maps}\\
A Pluto notebook, as with \ui{Pluto} & The cell's output & The cell's output, MapLibre GL JS loaded from jsDelivr at
  the version Pollis ships\\
\bottomrule
\end{tabularx}
\end{center}
The \ui{Plots} tab also draws pictures of other kinds that a Julia package shows: PNG, SVG and GIF images, Vega and
Vega-Lite charts, and Plotly charts.

\section{When ECharts is not enough}\label{sec:beyond}
Some pictures need a library made for them, and the panels that use one open a page of their own in Pollis's browser.
\begin{itemize}
  \item \textbf{Maps}, drawn by MapLibre GL JS (\S\ref{sec:maps}).
  \item \textbf{Street-level imagery}, the panels of Mapillary in Smart Regions: MapillaryJS, through
    \file{Pollis.streetview}.
  \item \textbf{Photorealistic 3D cities}, the Google Maps panel of Smart Regions: the Maps JavaScript API with its 3D
    maps, with the user's own key. Its terms allow the 3D tiles to be looked at only, not analysed, extracted or kept
    offline.
  \item \textbf{Point clouds, 3D city models and large layers}: deck.gl is the library meant for them; Pollis does not
    draw them yet.
  \item \textbf{3D charts}, such as surfaces and 3D scatter plots: echarts-gl, an extension of ECharts, which Pollis
    does not ship yet.
\end{itemize}

\section{Still to come}\label{sec:next}
\begin{itemize}
  \item \textbf{PollisECharts.jl}, written and registered, and the panels moved to it from ECharts.jl.
  \item \textbf{The panels that still draw in Julia}: on 21 November 2026, 168 panels of Pollis and its toolboxes
    draw with Plots.jl or StatsPlots.jl, and 3 with Makie.jl. Each is to move to ECharts.
  \item \textbf{What ECharts does not draw}: contour plots with labelled levels, mathematics in labels, and still images,
    SVG or PDF, for the reports Pollis writes in LaTeX.
  \item \textbf{3D}: deck.gl for point clouds and city models, echarts-gl for 3D charts.
\end{itemize}

\section{Where things are}
\begin{center}\small
\renewcommand{\arraystretch}{1.2}
\begin{tabularx}{\linewidth}{@{}p{0.5\linewidth}L@{}}
\toprule
\textbf{File} & \textbf{Role}\\
\midrule
\file{extensions/language-julia/src/interactive/plots.ts} & The \ui{Plots} tab: each kind of picture it draws, and
  ECharts.jl's charts drawn with the tab's own copy of ECharts\\
\file{extensions/language-julia/libs/echarts} & ECharts 6.0.0\\
\file{extensions/language-julia/libs/maplibre}, \file{.../libs/mapillary} & MapLibre GL JS 5.24.0, MapillaryJS 4.1.2\\
\file{extensions/language-julia/scripts/terminalserver/Pollis.jl} & \file{Pollis.webmap}, \file{Pollis.layer},
  \file{Pollis.tiles}, \file{Pollis.streetview}, and the pages of the maps\\
\file{extensions/language-julia/src/interactive/pollisMap.ts} & The drawing of a map, shared by its page and notebook
  cells\\
\file{extensions/language-julia/src/interactive/pollisMapPage.ts} & The page of a map in Pollis's browser\\
\file{extensions/language-julia/package.json} & The commands of the \ui{Plots} tab and the renderer \ui{Pollis Maps}\\
\bottomrule
\end{tabularx}
\end{center}

\references
\dref{CSRBB2023a}{Christ, S., Schwabeneder, D., Rackauckas, C., Borregaard, M.\,K. and Breloff, T. (2023). Plots.jl -- A User Extendable Plotting API for the Julia Programming Language. \emph{Journal of Open Research Software} 11, 5. \ddoi{10.5334/jors.431}}
\dref{DK2021a}{Danisch, S. and Krumbiegel, J. (2021). Makie.jl: Flexible high-performance data visualization for Julia. \emph{Journal of Open Source Software} 6(65), 3349. \ddoi{10.21105/joss.03349}}
\dref{Ero2026a}{Ero, Z. (2026). OpenFreeMap: Free and Open-Source Map Hosting with OpenStreetMap Data. \url{https://github.com/hyperknot/openfreemap}}
\dref{Goo2026a}{Google (2026a). Maps JavaScript API Documentation. \url{https://developers.google.com/maps/documentation/javascript}}
\dref{Goo2026b}{Google (2026b). 3D Maps in Maps JavaScript. \url{https://developers.google.com/maps/documentation/javascript/3d/overview}}
\dref{LMSSZ...2018a}{Li, D., Mei, H., Shen, Y., Su, S., Zhang, W., Wang, J., Zu, M. and Chen, W. (2018). ECharts: A Declarative Framework for Rapid Construction of Web-Based Visualization. \emph{Visual Informatics} 2(2), 136--146. \ddoi{10.1016/j.visinf.2018.04.011}}
\dref{MLC2026a}{MapLibre Contributors (2026). MapLibre GL JS: Interactive Vector Tile Maps in the Browser. \url{https://github.com/maplibre/maplibre-gl-js}}
\dref{Map2023a}{Mapillary (2023). MapillaryJS: Interactive, Extendable Street Imagery Map Experiences in the Browser. \url{https://github.com/mapillary/mapillary-js}}
\dref{Zwi2026a}{Zwitch, R. (2026). ECharts.jl: Julia Package for the Apache ECharts v6 Visualization Library. \url{https://github.com/randyzwitch/ECharts.jl}}

\end{document}
